\documentclass[10pt,xcolor=pdflatex,hyperref={unicode}]{beamer}
\usepackage{newcent}
\usepackage[utf8]{inputenc}
\usepackage[czech]{babel}
\usepackage[T1]{fontenc}
\usepackage{hyperref}
\usepackage{fancyvrb}
\usepackage{listings}
\usetheme{FIT}
\usepackage{graphicx}

%%%%%%%%%%%%%%%%%%%%%%%%%%%%%%%%%%%%%%%%%%%%%%%%%%%%%%%%%%%%%%%%%%
\title[Qick sort]{Řadící algoritmy: Quick sort}

\author[]{Štěpán Czajkowski}

%\institute[]{Brno University of Technology, Faculty of %Information Technology\\
%Bo\v{z}et\v{e}chova 1/2. 612 66 Brno - Kr\'alovo Pole\\
%login@fit.vutbr.cz}

\institute[]{Fakulta informačních technologií
Vysokého učení technického v Brně\\
Bo\v{z}et\v{e}chova 1/2. 612 66 Brno - Kr\'alovo Pole\\
xczajk01@stud.fit.vutbr.cz}

%České logo - Czech logo
%beamerouterthemeFIT.sty řádek 9: fitlogo1_cz

\date{30 duben, 2020}
%\date{\today}
%\date{} % bez data / without date

%%%%%%%%%%%%%%%%%%%%%%%%%%%%%%%%%%%%%%%%%%%%%%%%%%%%%%%%%%%%%%%%%%

\begin{document}

\frame[plain]{\titlepage}

\begin{frame}\frametitle{Motivace}
    \begin{itemize}
        \item Získání bodů do předmětu ITY 
        \item Zopakovat si Quicksort
        \item Naučit se prezentace v \LaTeX   \end{itemize}
    %Example \emph{content}.
\end{frame}


\begin{frame}\frametitle{Historie}
    \begin{itemize}
        \item Vytvořil Tony Hoare 
        \item Návrh 1959 
        \item Publikace 1961 v časopise Communications of the ACM
    \end{itemize}
    \begin{figure}[h]
        \includegraphics[width=7cm, height=5cm]{Bubble_original_latex.png}
    \end{figure}
\end{frame}

\begin{frame}\frametitle{Pivot}
    \begin{itemize}
        \item Prvek se kterým se porovnává zbytek pole
        \item Ideálně medián hodnot
        \item Různe metody náhodně, první nebo poslední prvek
        \item Metoda 3 mediánů 
        \begin{enumerate}
            \item vezmeme první, prostřední a poslední prvek 
            \item největší hodnotu na začátek pole nejmenší na konec pole prostřední hodnot pivot
        \end{enumerate}
        \item Dá se použít i na více hodnot
        
    \end{itemize}
\end{frame}

\begin{frame}\frametitle{Jak to funguje}
\begin{enumerate}
    \item<1-> Vyberte pivot a přesuntě ho na konec pole 
    \item<2-> Najděte první prvek z leva, větší než pivot 
    \item<3-> Najděte první prvek z prava, menší než pivot
    \item<4-> Pokud je pravý prvek víc v pravo než levý prvek, prvky vyměníme
    \item<5-> Pokud je prvek z leva víc v pravo než prvek z prava vyměníme prvek z leva za pivot
    \item<6-> pivot rozdělil pole na dvě části opakujeme na obou částech bez původního pivotu
\end{enumerate}
\end{frame}

\begin{frame}{Jak to funguje graficky}
    \includegraphics[width=9cm, height=8cm]{bubblesort.png}
\end{frame}

\begin{frame}{Pseudokód}
    \includegraphics[width=8cm, height=5cm]{pseudocode_quick.png}
\end{frame}

\begin{frame}\frametitle{Složitost}
    \begin{itemize}
        \item Nejhorší  $O(n^2)$\\ pokud je pivot největší nebo nejmenší prvek v poli
        \item Průměrná  $O(n\cdot \log_n)$
    \end{itemize}
        \begin{figure}
            \includegraphics[width=7cm, height=5cm]{quicksort_complexity.png}
        \end{figure}

\end{frame}

\bluepage{Děkuji za pozornost !}

\begin{frame}\frametitle{zdroje}
    \url{https://cs.wikipedia.org/wiki/Rychl\%C3\%A9_\%C5\%99azen\%C3\%AD}
    \url{https://www.youtube.com/watch?v=Hoixgm4-P4M}
    \url{https://www.javatpoint.com/quick-sort}
    \url{https://en.wikipedia.org/wiki/Quicksort}
    \url{https://cacm.acm.org/magazines/1961/7/14472-algorithm-64-quicksort/abstract}
    \url{https://www.interviewbit.com/tutorial/quicksort-algorithm/}
\end{frame}

\end{document}
